\documentclass[10pt,a4paper]{amsart}
\usepackage[margin=19mm]{geometry}
\usepackage{amsmath,amssymb,mathtools}
\usepackage[scr=rsfs,cal=euler]{mathalfa}
\usepackage{xfrac}
\usepackage[unicode,hidelinks,bookmarks=false]{hyperref}
\newcommand{\ie}{i.e.\ }
\newcommand{\cf}{cf.\ }
\newcommand{\resp}{respectively\ }
\newcommand{\Subsection}{\subsection}
\providecommand{\nobreakdash}{\nobreak\hbox{-}}
\setlength{\emergencystretch}{2em}
\multlinegap0pt
\usepackage{bbm}
\usepackage{amssymb}
\usepackage{tikz}
\usepackage[shortlabels]{enumitem}
\newtheorem{theorem}{Theorem}
\newtheorem{proposition}{Proposition}
\newtheorem{lemma}{Lemma}
\newcommand{\ud}{\mathrm{d}}
\title{Existence and nonexistence of normalized solutions for nonlinear Schr\"{o}dinger equation involving combined nonlinearities in bounded domain}
\author{Zhen-Feng Jin, Weimin Zhang}
\date{Source: arXiv:2602.16263v1 (CC BY 4.0)}
\begin{document}
\maketitle

\noindent\emph{Source and licence.} Existence and nonexistence of normalized solutions for nonlinear Schr\"{o}dinger equation involving combined nonlinearities in bounded domain. Version arXiv:2602.16263v1 (CC BY 4.0), distributed by arXiv under the Creative Commons Attribution 4.0 International licence, http://creativecommons.org/licenses/by/4.0/, as recorded in the arXiv licence metadata for this exact version and shown on the abstract page https://arxiv.org/abs/2602.16263v1. Full licence notice: \texttt{licenses/arxiv-2602.16263-CC-BY-4.0.txt}.\par\smallskip

\noindent\textit{Editorial scope.} Counted unit: the article's theorem T041801 (source0.tex lines 344--347). The article's complete original proof of it is delivered with it (source0.tex lines 579--651, one \texttt{proof} environment of 73 lines, which the article places away from the statement). No mathematics is written, repaired or re-derived: the statement and the proof are the article's own lines, in the article's own notation, and no retained line is substituted or re-flowed.  Retained uncounted as the source-local context the proof needs: lines 211--219; lines 331--335; lines 337--342; lines 461--465; lines 476--558; lines 559--578.  Nothing was omitted from any retained span, so the package carries no omission record.  No retained line contains a citation, so the wrapper carries no bibliography and no bibliography entry.  No retained line contains a figure environment, an image-inclusion command or drawing code, so no image is delivered and the package declares no asset dependency.  Editorial formatting only, in addition: an AMS \texttt{amsart} preamble that restates the article's own declarations and macros used by the retained lines, namely the environments \texttt{theorem}, \texttt{proposition}, \texttt{lemma} on separate counters, together with the standard packages those lines need.  The retained environments print in this document as Theorem 1, Proposition 1, Lemma 1 in order of appearance, whereas the article prints the same items with the numbering of its own full document.  The complete native source of the article as distributed in the arXiv e-print for this exact version is bundled unchanged as \texttt{sources/194864/source0.tex}.
This paper is concerned with the existence and nonexistence of solutions for the following nonlinear Schr\"{o}dinger equation involving combined power nonlinearities \begin{equation}\label{041003}
\begin{cases}
\begin{aligned}
&-\Delta u+\omega u=\mu u^{p-1}+u^{q-1},~ u>0 \quad &&\text { in } \Omega, \\
&u=0  &&\text { on } \partial\Omega, \\
&\|u\|_2^2=\rho>0,
\end{aligned}
\end{cases}
\end{equation}

\begin{equation}\label{2505261720}
\begin{aligned}
m_{\alpha}:=\inf_{A_{\alpha}} I(u),\quad \widetilde{m}_{\alpha}:=\inf_{\partial A_{\alpha}} I(u).
\end{aligned}
\end{equation}

\begin{equation}\label{2505271451}
\rho<
\begin{aligned}
\frac{1}{\alpha-\lambda_1(\Omega)}\Big(\frac{2^{*}}{2}\mathcal{S}^{\frac{2^{*}}{2}}\Big)^{\frac{2}{2^{*}-2}}.
\end{aligned}
\end{equation}

\begin{equation}\label{050501}
\begin{aligned}
\mathcal{S}\|u\|^{2}_{L^{2^{*}}(\mathbb{R}^N)} \leq\|\nabla u\|^{2}_{L^{2}(\mathbb{R}^N)}\quad \forall\, u \in D^{1,2}(\mathbb{R}^{N}), \quad \text{ (Sobolev inequality) }
\end{aligned}
\end{equation}

\begin{proposition}\label{P050601}
Let $N\geq 3$, $2<p<q=2^*$, $\mu\in \mathbb{R}$ and $\alpha>\lambda_{1}(\Omega)$. If \eqref{2505271451} holds, then any minimizing sequence associated to $m_{\alpha}$ {\rm(}given in \eqref{2505261720}{\rm)} is relatively compact in $A_{\alpha}$. In particular, $m_{\alpha}$ can be achieved.
\begin{proof}
Let $\{u_n\}\subset H^1_0(\Omega)$ be a minimizing sequence for $m_{\alpha}$. That is,
\begin{equation}\label{040503}
\begin{cases}
\|u_n\|^{2}_{L^{2}(\Omega)}=\rho, \\
\|\nabla u_n\|^{2}_{L^{2}(\Omega)}\leq\rho\alpha, \\
m_{\alpha}\leq I(u_n)\leq m_{\alpha}+o(1)
\end{cases}
\end{equation}
as $n \rightarrow \infty$. Up to a subsequence, there exists $u_0\in H^1_0(\Omega)$ such that
\begin{equation*}
\begin{cases}
u_n\rightharpoonup u_0 \text{ in } H^1_0(\Omega), \\
\|\nabla u_0\|^{2}_{L^{2}(\Omega)}\leq\liminf\limits_{n\rightarrow \infty}\|\nabla u_n\|^{2}_{L^{2}(\Omega)}, \\
\|u_0\|^{2}_{L^{2}(\Omega)}=\rho.
\end{cases}
\end{equation*}
So, $u_0\in A_{\alpha}$ and $I(u_0)\geq m_{\alpha}$.

Let $\tilde{u}_n=u_n-u_0$. Then, we extract a subsequence, still denoted by $\{\tilde{u}_n\}$, such that
\begin{equation}\label{2505211318}
\begin{aligned}
\tilde{u}_n&\rightharpoonup 0 ~~\text{ in } H^1_0(\Omega), \\
\tilde{u}_n&\rightarrow 0 ~~\text{ in } L^{r}(\Omega) \text{ for } r\in[1,2^{*}),\\
\tilde{u}_n&\rightarrow 0 ~~\text{ a.e.}\text{ in }\Omega
\end{aligned}
\end{equation}
as $n \rightarrow \infty$. By weak convergence and Br\'{e}zis-Lieb lemma, we deduce that
\begin{equation}\label{040504}
\begin{aligned}
\|\nabla (\tilde{u}_n+u_0)\|^2_{L^{2}(\Omega)}
&=\|\nabla \tilde{u}_n\|^2_{L^{2}(\Omega)}+\|\nabla u_0\|^2_{L^{2}(\Omega)}+o(1), \\
\|\tilde{u}_n+u_0\|^{2^{*}}_{L^{2^{*}}(\Omega)}
&=\|\tilde{u}_n\|^{2^{*}}_{L^{2^{*}}(\Omega)}+\|u_0\|^{2^{*}}_{L^{2^{*}}(\Omega)}+o(1).
\end{aligned}
\end{equation}
By \eqref{040504}, we get
\begin{equation*}
\begin{aligned}
I(u_n)&=\frac{1}{2}\int_{\Omega}|\nabla u_n|^2 \ud x
-\frac{\mu}{p}\int_{\Omega}|u_n|^{p} \ud x
-\frac{1}{2^{*}}\int_{\Omega}|u_n|^{2^{*}} \ud x\\
&=I(\tilde{u}_n)+I(u_0)+o(1).
\end{aligned}
\end{equation*}
Using \eqref{040503} and $I(u_0)\geq m_{\alpha}$, we have $I(\tilde{u}_n)\leq o(1)$ as $n \rightarrow \infty$. From \eqref{050501} and \eqref{2505211318}, it follows that
\begin{equation}\label{050604}
\begin{aligned}
\int_{\Omega}|\nabla \tilde{u}_n|^2 \ud x
&\leq\frac{2\mu}{p}\int_{\Omega}|\tilde{u}_n|^{p} \ud x
+\frac{2}{2^{*}}\int_{\Omega}|\tilde{u}_n|^{2^{*}} \ud x+o(1)\\
&\le\frac{2}{2^{*}}\mathcal{S}^{-\frac{2^{*}}{2}}\|\nabla \tilde{u}_n\|^{2^{*}}_{L^{2}(\Omega)}+o(1).
\end{aligned}
\end{equation}
Since $\{\|\nabla \tilde{u}_n\|^{2}_{L^{2}(\Omega)}\}$ is a bounded sequence, up to a subsequence, we distinguish the two cases
\begin{equation*}
\begin{aligned}
\text{either }\mathrm{(a)}~~\|\nabla \tilde{u}_n\|^{2}_{L^{2}(\Omega)}\rightarrow 0 \quad
\text{ or } \mathrm{(b)}~~\|\nabla \tilde{u}_n\|^{2}_{L^{2}(\Omega)}\rightarrow \ell>0.
\end{aligned}
\end{equation*}
If $\mathrm{(b)}$ holds,
%\begin{equation*}
%\begin{aligned}
%\|\nabla \tilde{u}_n\|^{2}_{L^{2}(\Omega)}
%\geq\Big(\frac{2^{*}}{2}\mathcal{S}^{\frac{2^{*}}{2}}\Big)^{\frac{2}{2^{*}-2}}+o(1).
%\end{aligned}
%\end{equation*}
by \eqref{040503}-\eqref{050604} and $\|\nabla u_0\|^{2}_{L^{2}(\Omega)}\geq\lambda_{1}(\Omega)\rho$ with $u_0\in S_{\rho}$, we get that
\begin{equation*}
\begin{aligned}
\Big(\frac{2^{*}}{2}\mathcal{S}^{\frac{2^{*}}{2}}\Big)^{\frac{2}{2^{*}-2}}
\leq\|\nabla \tilde{u}_n\|^{2}_{L^{2}(\Omega)}+o(1)
&=\|\nabla u_n\|^{2}_{L^{2}(\Omega)}
-\|\nabla u_0\|^{2}_{L^{2}(\Omega)}+o(1)\\
&\leq(\alpha-\lambda_{1}(\Omega))\rho+o(1),
\end{aligned}
\end{equation*}
which contradicts \eqref{2505271451}. Thus, $\|\nabla \tilde{u}_n\|^{2}_{L^{2}(\Omega)}\rightarrow 0$. That is, $u_n\rightarrow u_0$ in $H^1_0(\Omega)$ and $I(u_0)=m_{\alpha}$.
\end{proof}
\end{proposition}

\begin{lemma}\label{L050602}
Let $2<p<q\le 2^*$, $\mu\in \mathbb{R}$. Assume that there exist $\alpha_2>\alpha_1\ge\lambda_1(\Omega)$ such that 
\begin{equation}\label{050603}
\begin{aligned}
 \widetilde{m}_{\alpha_1}<\widetilde{m}_{\alpha_2}.
\end{aligned}
\end{equation}
If $N\ge 3$ and $q=2^*$, we need the extra assumption that \eqref{2505271451} holds with $\alpha= \alpha_2$. Then $m_{\alpha_2}<\widetilde{m}_{\alpha_2}$, and $m_{\alpha_2}$ can be achieved by a positive solution of problem \eqref{041003}.
\begin{proof}
By \eqref{050603}, we have
\begin{equation}\label{2510122022}
\begin{aligned}
m_{\alpha_2} = \inf\{\widetilde{m}_{\alpha}:\lambda_1(\Omega)\leq \alpha \leq \alpha_2\}
\leq \widetilde{m}_{\alpha_1}
<\widetilde{m}_{\alpha_2}.
\end{aligned}
\end{equation}
When $q<2^*$, since  the embeddings $H_0^1(\Omega)\subset L^p(\Omega)$ and $H_0^1(\Omega)\subset L^q(\Omega)$ are compact, it is clear that there exists some $u_0\in A_{\alpha_2}$ such that $I(u_0)=m_{\alpha_2}$. When $q=2^*$, from Proposition \ref{P050601}, we know that $m_{\alpha_2}$ can  be also achieved by some $u_0\in A_{\alpha_2}$. Next, it follows from \eqref{2510122022} that $u_0\in A_{\alpha_2}\setminus \partial A_{\alpha_2}$.
\end{proof}
\end{lemma}

\begin{theorem}\label{T041801}
Let $\mu\in \mathbb{R}$ and $2<p\le q\le  2^*$. Then there exists a $\rho_0=\rho_0(N,p,q,\Omega,\mu)>0$ such that
 for any $\rho\in(0,\rho_0)$,  \eqref{041003} has a solution $(\omega, u)\in\mathbb{R}\times S_\rho$, which is a local minimizer of $I|_{S_\rho}$.
\end{theorem}

\begin{proof}[\bf Proof of Theorem \ref{T041801}]
By $\sqrt{\rho}\varphi_1 \in \partial A_{\lambda_1(\Omega)}$,  then we have
%$\widetilde{m}_{\lambda_1(\Omega)}\leq \frac{\rho}{2}\lambda_1(\Omega)$.
\begin{equation}\label{2505211445}
\begin{aligned}
\widetilde{m}_{\lambda_1(\Omega)}\leq I(\sqrt{\rho}\varphi_1)
=\frac{\lambda_1(\Omega)}{2}\rho
-\frac{\mu}{p}\|\varphi_1\|^p_{L^{p}(\Omega)}\rho^{\frac{p}{2}}
-\frac{1}{2^{*}}\|\varphi_1\|^{q}_{L^{q}(\Omega)}\rho^{\frac{q}{2}}.
\end{aligned}
\end{equation}
This proof will be split into three cases.
\smallskip

{\it Case $N\ge 3$, $2<p<q=2^*$}. For $u\in \partial A_{\alpha}$, by Sobolev inequality and Gagliardo-Nirenberg inequality, it holds that
\begin{equation}\label{25052719099}
\begin{aligned}
I(u)&=\frac{1}{2}\int_{\Omega}|\nabla u|^2 \ud x
-\frac{\mu}{p}\int_{\Omega}|u|^{p} \ud x
-\frac{1}{2^{*}}\int_{\Omega}|u|^{2^{*}} \ud x\\
&\geq \frac{1}{2}\int_{\Omega}|\nabla u|^2 \ud x
-\frac{|\mu|}{p}\mathcal{C}^{p}_{N,p}\|\nabla u\|^{p\gamma_{p}}_{L^{2}(\Omega)}
\|u\|^{p(1-\gamma_{p})}_{L^{2}(\Omega)}
-\frac{1}{2^{*}}\mathcal{S}^{-\frac{2^{*}}{2}}\|\nabla u\|^{2^{*}}_{L^{2}(\Omega)}\\
&=\frac{\alpha}{2}\rho
-\frac{|\mu|}{p}\mathcal{C}^{p}_{N,p}\alpha^{\frac{p\gamma_{p}}{2}}\rho^{\frac{p}{2}}
-\frac{1}{2^{*}}\mathcal{S}^{-\frac{2^{*}}{2}}\alpha^{\frac{2^{*}}{2}}\rho^{\frac{2^{*}}{2}}.
\end{aligned}
\end{equation}
Thus, we get that
\begin{equation}\label{2505211446}
\begin{aligned}
\widetilde{m}_{\alpha}\geq
\frac{\alpha}{2}\rho
-\frac{|\mu|}{p}\mathcal{C}^{p}_{N,p}\alpha^{\frac{p\gamma_{p}}{2}}\rho^{\frac{p}{2}}
-\frac{1}{2^{*}}\mathcal{S}^{-\frac{2^{*}}{2}}\alpha^{\frac{2^{*}}{2}}\rho^{\frac{2^{*}}{2}}.
\end{aligned}
\end{equation}

For each $\alpha>\lambda_1(\Omega)>0$, we consider the function
\begin{equation*}
\begin{aligned}
\hat{f}_{\alpha}(\rho)
=&\frac{\alpha-\lambda_1(\Omega)}{2}\rho
+{\frac{1}{p}\left(\mu\|\varphi_1\|^p_{L^{p}(\Omega)}-|\mu|\mathcal{C}^{p}_{N,p}\alpha^{\frac{p\gamma_{p}}{2}}\right)\rho^{\frac{p}{2}}}
+\frac{1}{2^{*}}\big(\|\varphi_1\|^{2^{*}}_{L^{2^{*}}(\Omega)}
-\mathcal{S}^{-\frac{2^{*}}{2}}\alpha^{\frac{2^{*}}{2}}\big)\rho^{\frac{2^{*}}{2}}.\\
\end{aligned}
\end{equation*}
By $\alpha>\lambda_1(\Omega)$ and $2<p<2^*$, we deduce that there exists $\rho_0=\rho_0(N,p,\Omega, \mu, \alpha)>0$ such that  for any $\rho\in(0,\rho_0)$, $\hat{f}_{\alpha}(\rho)>0$.
Thus, it follows from \eqref{2505211445}, \eqref{2505211446} that
%Thus, for $\rho>0$ sufficiently small, we have
\begin{equation*}
\begin{aligned}
\widetilde{m}_{\alpha}
&\geq
\frac{\alpha}{2}\rho
-\frac{|\mu|}{p}\mathcal{C}^{p}_{N,p}\alpha^{\frac{p\gamma_{p}}{2}}\rho^{\frac{p}{2}}
-\frac{1}{2^{*}}\mathcal{S}^{-\frac{2^{*}}{2}}\alpha^{\frac{2^{*}}{2}}\rho^{\frac{2^{*}}{2}}\\
&>\frac{\lambda_1(\Omega)}{2}\rho
-\frac{\mu}{p}\|\varphi_1\|^p_{L^{p}(\Omega)}\rho^{\frac{p}{2}}
-\frac{1}{2^{*}}\|\varphi_1\|^{2^{*}}_{L^{2^{*}}(\Omega)}\rho^{\frac{2^{*}}{2}}
\geq\widetilde{m}_{\lambda_1(\Omega)},\quad \forall\,\rho\in(0,\rho_0).
\end{aligned}
\end{equation*}
Using Lemma \ref{L050602}, we get that $m_{\alpha}$ is achieved by a positive solution of problem \eqref{041003}.

\smallskip
{\it Case $N\ge 1$, $2<p<q<2^*$}. In this case,  the Sobolev inequality can be replaced by Gagliardo-Nirenberg inequality, all steps above can proceed.

\smallskip
{\it Case $N\ge 1$, $2<p=q\le 2^*$}. If $\mu>-1$,  this situation is identical to the above cases with $\mu=0$. If $\mu\le -1$, the associated energy is coercive, and it is simple to obtain a global minimizer. Moreover, the global minimizer exists for all $\rho>0$.
\end{proof}

\end{document}
